\documentclass[11pt,aspectratio=169]{beamer}

\usetheme{metropolis}
\usecolortheme{default}
\metroset{numbering=fraction, progressbar=frametitle, sectionpage=none}

\setbeamertemplate{navigation symbols}{}
\setbeamercolor{frametitle}{bg=mDarkTeal, fg=white}
\setbeamercolor{progress bar}{fg=mDarkTeal, bg=mDarkTeal!20}
\definecolor{mDarkTeal}{HTML}{23373B}
\definecolor{mAccent}{HTML}{EB811B}

\usepackage{amsmath,amssymb}
\usepackage{booktabs}

\title{Math Bootcamp}
\subtitle{Day 5 --- Dynamic Programming (Basics)}
\author{Math Camp 2026}
\date{}

\newcommand{\QFrame}[1]{\begin{frame}{#1 \quad {\small\textcolor{mAccent}{Together}}}}
\newcommand{\THFrame}[1]{\begin{frame}{#1 \quad {\small\textcolor{mAccent}{Take home}}}}
\newcommand{\InterestOnly}{\textcolor{mAccent}{\textbf{for interest only}}}

\begin{document}

\begin{frame}
  \titlepage
\end{frame}

\begin{frame}{Day 5 roadmap}
  \small
  \begin{enumerate}
    \item \textbf{5.1} Sequential decisions and the DP idea
    \item \textbf{5.2} State, control, and law of motion
    \item \textbf{5.3} Two-period consumption--savings model
    \item \textbf{5.4} Backward induction
    \item \textbf{5.5} Bellman equation and policy function
    \item \textbf{5.6} Euler equation and link to standard optimization
  \end{enumerate}

  \medskip
  \footnotesize
  \textbf{Focus today:} one worked \textbf{two-period} example (finite horizon).
\end{frame}

%==============================================================================
\section{5.1 --- Sequential Decisions}
%==============================================================================

\begin{frame}{5.1 Why dynamic programming?}
  \small
  Many economic problems are \textbf{sequential}:
  \begin{itemize}
    \item You choose something \textbf{today}.
    \item It changes what you \textbf{can} choose tomorrow.
    \item Today's choice trades off \textbf{now} vs.\ \textbf{later}.
  \end{itemize}

  \medskip
  \textbf{Examples.}\
  Consumption vs.\ savings; hiring vs.\ waiting; extracting a resource vs.\ leaving it in the ground.

  \medskip
  \textbf{Dynamic programming (DP).}\
  Break a multi-period problem into \textbf{smaller subproblems}, linked by a \textbf{value function}.
\end{frame}

\begin{frame}{5.1 Finite horizon and the Markov idea}
  \small
  \textbf{Finite horizon.}\
  Time $t=1,2,\ldots,T$ (today: $T=2$).\
  After period $T$, the problem ends (\textbf{terminal condition}).

  \medskip
  \textbf{Markov structure (informal).}\
  For decision-making, \textbf{only the current state matters} --- not the full history, once the state is given.\
  Example: tomorrow's opportunities depend on \textbf{assets on hand}, not on every past consumption choice separately.

  \medskip
  \textbf{Key objects.}\
  \textbf{State} (where you are), \textbf{control} (what you choose), \textbf{transition} (how state moves).
\end{frame}

%==============================================================================
\section{5.2 --- State, Control, Transition}
%==============================================================================

\begin{frame}{5.2 State versus control}
  \small
  \begin{tabular}{@{}ll@{}}
    \toprule
    \textbf{Object} & \textbf{Meaning} \\
    \midrule
    \textbf{State} $x_t$ & Variables summarizing the situation at the \textbf{start} of period $t$ \\
    \textbf{Control} $u_t$ & Variables chosen \textbf{in} period $t$ \\
    \textbf{Transition} & Law of motion $x_{t+1}=g(x_t,u_t)$ \\
    \bottomrule
  \end{tabular}

  \medskip
  \textbf{Feasibility.}\
  Controls must satisfy constraints each period (budget, nonnegativity, etc.).

  \medskip
  \textbf{Payoff.}\
  Period-$t$ utility $u(c_t)$ (or profit, welfare, \ldots).\
  Future payoffs are discounted by $\beta\in(0,1]$.
\end{frame}

%==============================================================================
\section{5.3 --- Two-Period Model}
%==============================================================================

\begin{frame}{5.3 Two-period consumption--savings}
  \small
  \textbf{Story.}\
  A household lives for \textbf{two periods}, $t=1,2$.\
  Income $y>0$ in period~1 only; no labor income in period~2.\
  It chooses consumption $c_1,c_2\ge 0$ and saves
  \[
    s=y-c_1 \qquad (\text{savings in period 1}).
  \]
  Savings earn gross return $1+r$ ($r>-1$).\
  Period-2 consumption is
  \[
    c_2=(1+r)s=(1+r)(y-c_1).
  \]

  \medskip
  \textbf{Preferences.}\
  \[
    \max_{c_1,c_2}\ \log(c_1)+\beta\log(c_2),
    \qquad \beta\in(0,1).
  \]
  Log utility keeps algebra clean; same logic works for $u(\cdot)$ more generally.
\end{frame}

\begin{frame}{5.3 State and control in this model}
  \small
  \textbf{State at start of period $t$:} \textbf{assets} $b_t$ (wealth before choosing $c_t$).

  \medskip
  \begin{itemize}
    \item \textbf{Period 1:} $b_1=0$.\
      Income $y$ arrives; choose $c_1$.\
      Assets carried to period~2:
      \[
        b_2=(1+r)(b_1+y-c_1)=(1+r)(y-c_1).
      \]
    \item \textbf{Period 2:} state is $b_2$.\
      Choose $c_2$ subject to $0\le c_2\le b_2$ (consume all assets by the end).
  \end{itemize}

  \medskip
  \textbf{Controls:} $u_1=c_1$, $u_2=c_2$.\
  \textbf{Transition:} $b_{t+1}=(1+r)(b_t+y_t-c_t)$ with $y_2=0$.
\end{frame}

%==============================================================================
\section{5.4 --- Backward Induction}
%==============================================================================

\begin{frame}{5.4 Backward induction --- idea}
  \small
  \textbf{Strategy.}\
  In a \textbf{finite-horizon} problem, solve from the \textbf{last period backward}:
  \begin{enumerate}
    \item Period~2: optimize taking the state $b_2$ as \textbf{given}.
    \item Period~1: optimize knowing \textbf{optimal future behavior} is already solved.
  \end{enumerate}

  \medskip
  \textbf{Why backward?}\
  In period~2 there is \textbf{no future left} --- only current utility matters.\
  Period~1 must account for how today's savings affect tomorrow's feasible consumption.
\end{frame}

\begin{frame}{5.4 Step 1 --- last period ($t=2$)}
  \small
  \textbf{Value function in period 2} (state $b_2$):
  \[
    V_2(b_2)=\max_{0\le c_2\le b_2}\ \log(c_2).
  \]

  \medskip
  \textbf{Solution.}\
  Strictly increasing $\log(\cdot)$ $\Rightarrow$ consume everything:
  \[
    c_2^*(b_2)=b_2,
    \qquad
    V_2(b_2)=\log(b_2).
  \]

  \medskip
  \textbf{Policy function} $c_2^*(b_2)$: maps state $\to$ optimal control.\
  Here: ``use all assets.''
\end{frame}

\begin{frame}{5.4 Step 2 --- period $t=1$}
  \small
  \textbf{State:} $b_1=0$, income $y$.\
  Choosing $c_1$ sets $b_2=(1+r)(y-c_1)$.

  \medskip
  \textbf{Problem:}
  \[
    V_1(0)=\max_{0<c_1<y}\ \Bigl\{\log(c_1)+\beta V_2\bigl((1+r)(y-c_1)\bigr)\Bigr\}.
  \]
  Substitute $V_2(b_2)=\log(b_2)$:
  \[
    V_1(0)=\max_{0<c_1<y}\ \Bigl\{\log(c_1)+\beta\log\bigl((1+r)(y-c_1)\bigr)\Bigr\}.
  \]

  \medskip
  \textbf{FOC} ($u=\log$):
  \[
    \frac{1}{c_1}=\frac{\beta}{y-c_1}
    \quad\Longrightarrow\quad
    c_1^*=\frac{y}{1+\beta}.
  \]
  Then $c_2^*=(1+r)(y-c_1^*)=\dfrac{\beta(1+r)y}{1+\beta}$.
\end{frame}

%==============================================================================
\section{5.5 --- Bellman Equation}
%==============================================================================

\begin{frame}{5.5 Bellman equation --- general form}
  \small
  \textbf{Principle of optimality.}\
  Optimal plan $\Rightarrow$ remaining decisions are optimal from whatever state remains.

  \medskip
  \textbf{Recursive value function} (finite $T$):
  \[
    V_T(x_T)=\max_{u_T\in\mathcal{F}_T(x_T)}\ u(x_T,u_T)
  \]
  \[
    V_t(x_t)=\max_{u_t\in\mathcal{F}_t(x_t)}\ \Bigl\{u(x_t,u_t)+\beta V_{t+1}\bigl(g(x_t,u_t)\bigr)\Bigr\},
    \quad t=1,\ldots,T-1.
  \]

  \medskip
  \textbf{Read the recursion.}\
  Current payoff $+$ discounted \textbf{continuation value} $V_{t+1}$ from next state.
\end{frame}

\begin{frame}{5.5 Bellman equation --- two-period special case}
  \small
  With $T=2$, $b_1=0$, and $V_2(b_2)=\log(b_2)$:
  \[
    V_1(0)=\max_{c_1}\ \Bigl\{\log(c_1)+\beta V_2\bigl((1+r)(y-c_1)\bigr)\Bigr\}.
  \]

  \medskip
  \textbf{Policy functions:}
  \[
    c_1^*(0)=\frac{y}{1+\beta},
    \qquad
    c_2^*(b_2)=b_2.
  \]

  \medskip
  \textbf{Value at start:}
  \[
    V_1(0)=\log\!\Bigl(\frac{y}{1+\beta}\Bigr)+\beta\log\!\Bigl(\frac{\beta(1+r)y}{1+\beta}\Bigr).
  \]
\end{frame}

%==============================================================================
\section{5.6 --- Euler Equation and Comparison}
%==============================================================================

\begin{frame}{5.6 Euler equation (two periods)}
  \small
  \textbf{Same optimum, different language.}\
  The FOC from backward induction is the \textbf{two-period Euler equation}:
  \[
    u'(c_1^*)=\beta(1+r)\,u'(c_2^*).
  \]
  For $\log$ utility, $u'(c)=1/c$, so
  \[
    \frac{1}{c_1^*}=\beta(1+r)\,\frac{1}{c_2^*}
    \quad\Longleftrightarrow\quad
    c_2^*=\beta(1+r)\,c_1^*.
  \]

  \medskip
  \textbf{Meaning.}\
  Marginal utility today equals discounted marginal utility tomorrow, times the return to saving.\
  If $r=0$: $c_2^*=\beta c_1^*$ (patient households consume more later only if $\beta<1$ is balanced by returns).
\end{frame}

\begin{frame}{5.6 Same answer from the lifetime budget constraint}
  \small
  \textbf{Combine constraints} ($b_1=0$):
  \[
    c_1+\frac{c_2}{1+r}=y.
  \]

  \medskip
  \textbf{Lagrangian:}
  \[
    \mathcal{L}=\log(c_1)+\beta\log(c_2)+\lambda\Bigl(y-c_1-\frac{c_2}{1+r}\Bigr).
  \]
  FOC give $1/c_1=\lambda$ and $\beta/c_2=\lambda/(1+r)$, hence $c_2^*=\beta(1+r)c_1^*$.\
  Substitute into the budget:
  \[
    c_1^*\Bigl(1+\beta\Bigr)=y
    \quad\Longrightarrow\quad
    c_1^*=\frac{y}{1+\beta}.
  \]

  \medskip
  \textbf{Takeaway.}\
  DP (backward) and static intertemporal optimization (one budget) are \textbf{two views of the same problem}.
\end{frame}

\begin{frame}{5.6 Beyond two periods \InterestOnly}
  \footnotesize
  \InterestOnly

  \medskip
  \textbf{More periods.}\
  Same backward logic: solve $V_T$, then $V_{T-1}$, \ldots, down to $V_1$.

  \medskip
  \textbf{Infinite horizon.}\
  Value function iteration: start with a guess $V^0$, apply the Bellman operator until convergence.\
  (Algorithm details: fall macro/micro courses.)

  \medskip
  \textbf{State may grow.}\
  With capital $k_{t+1}=(1-\delta)k_t+i_t$, the state is $k_t$; controls are $c_t,i_t$.\
  Same template: $V_t(k_t)=\max\{u(c_t)+\beta V_{t+1}(k_{t+1})\}$.
\end{frame}

\QFrame{5.6 Exercise: Two-period savings with CRRA}
  \small
  \textbf{Exercise.}\
  Same setup as above ($b_1=0$, income $y$ in period~1, return $1+r$), but
  $u(c)=\dfrac{c^{1-\sigma}}{1-\sigma}$ with $\sigma>0$, $\sigma\neq 1$.\
  Using backward induction, show
  \[
    c_1^*=\frac{y}{1+\beta^{1/\sigma}(1+r)^{(1-\sigma)/\sigma}},
    \qquad
    c_2^*=\beta^{1/\sigma}(1+r)^{1/\sigma}\,c_1^*.
  \]

  \textbf{Solution (sketch).}\
  $V_2(b_2)=\max_{c_2\le b_2} u(c_2)=u(b_2)$, so $c_2^*(b_2)=b_2$.\
  Period~1 FOC:
  \[
    c_1^{-\sigma}=\beta(1+r)^{1-\sigma}(y-c_1)^{-\sigma}
    \quad\Longrightarrow\quad
    \frac{c_2^*}{c_1^*}=\beta^{1/\sigma}(1+r)^{1/\sigma}.
  \]
  Plug into $c_1+c_2/(1+r)=y$ to get $c_1^*$.\
  Check: $\sigma=1$ (log) gives $c_1^*=y/(1+\beta)$.
\end{frame}

\THFrame{5.6 Exercise: Income in both periods}
  \small
  \textbf{Exercise.}\
  Income $y_1,y_2>0$ in periods 1 and 2; gross return $1+r$; utility $\log(c_1)+\beta\log(c_2)$.\
  State $b_1=0$; transition $b_2=(1+r)(y_1-c_1)$; period-2 budget $c_2\le b_2+y_2$.\
  Derive $V_2(b_2)$, the Bellman equation for $V_1(0)$, and the optimal $c_1^*,c_2^*$.

  \textbf{Solution.}\
  $V_2(b_2)=\max_{0<c_2\le b_2+y_2}\log(c_2)$ $\Rightarrow$ $c_2^*(b_2)=b_2+y_2$, $V_2(b_2)=\log(b_2+y_2)$.\
  \[
    V_1(0)=\max_{c_1}\ \Bigl\{\log(c_1)+\beta\log\bigl(y_2+(1+r)(y_1-c_1)\bigr)\Bigr\}.
  \]
  FOC: $1/c_1=\beta(1+r)/(y_2+(1+r)(y_1-c_1))$.\
  Lifetime budget: $c_1+c_2/(1+r)=y_1+y_2/(1+r)$ gives the same Euler equation
  $c_2^*=\beta(1+r)c_1^*$.
\end{frame}

\end{document}
